% --- LaTeX Essay Template - S. Venkatraman ---

% --- Set document class and font size ---

\documentclass[letterpaper, 11pt]{article}

% --- Package imports ---

% lipsum is just for generating placeholder text and can be removed
\usepackage{hyperref, lipsum} 

% --- Fonts (Set to Palatino or Times New Roman if desired) ---

% \usepackage{newpxtext, newpxmath}
% \usepackage{times}

% --- Page layout settiings ---

% Set page margins
\usepackage[left=1.2in, right=1.2in, top=.9in, bottom=.9in]{geometry}

% Anchor footnotes to the bottom of the page
\usepackage[bottom]{footmisc}
\usepackage{setspace}
\setstretch{1.15}
% Set line spacing
\renewcommand{\baselinestretch}{1.2}

% Use this for a one-off '*' footnote (e.g., a URL in a heading)
\newcommand{\starfootnote}[1]{%
\begingroup
\renewcommand{\thefootnote}{\fnsymbol{footnote}}%
\footnote{#1}%
\endgroup
\addtocounter{footnote}{-1}%
}

% Set spacing between paragraphs
\setlength{\parskip}{2mm}

% --- Page formatting settings ---



% --- Essay title and author ---
\title{\vspace{-1.5cm} Summary of \textit{Veridical Data Science}, Chapters 9 \footnote{\url{https://vdsbook.com/}}}

\date{\today}

% --- Main text ---
% --- Document starts here ---

\begin{document}
% Set link colors for labeled items and URLs (blue) 
\hypersetup{colorlinks=true, linkcolor=blue, urlcolor=blue}
\maketitle

\section*{Chapter 9: Continuous Responses \& Least Squares}

Chapter 9 introduces loss functions and predictive evaluation.
It compares least squares and least absolute deviation, showing how different loss functions prioritize errors differently.
Squared loss penalizes large errors more heavily, while absolute loss is more robust to outliers.
We should choose the loss function based on the context and the importance of different types of errors in our specific application.

The chapter explains metrics such as MSE, rMSE, and MAE, highlighting that model rankings can differ across measures.
It stresses that training error underestimates future error, making validation essential.
To prevent overfitting, we should use validation sets, consider model complexity, and make sure model stability is good.
% --- Document ends here ---
\end{document}
